
%%%%%%%%%%%%%%%%%%%%%%%%%%%%%%%%%%%%%%%%%%%%%%%%%%%%%%%%%%%%%%%%%%%%%%%%%%%%%%
%%%%%%%%%%%%%%%%%%%%%%%%%%%%%%%%%%%%%%%%%%%%%%%%%%%%%%%%%%%%%%%%%%%%%%%%%%%%%% 

%\subsection{Perfect Mixability of Near-Final Configurations}
%\label{subsec: near-final configurations}
%
%%%%%%%%%%%%%%%%%%%%%%%%%%%%%%%%%%%%%%%%%%%%%%%%%%%%%%%%%%%%%%%%%%%%%%%%%%%%%%
%%%%%%%%%%%%%%%%%%%%%%%%%%%%%%%%%%%%%%%%%%%%%%%%%%%%%%%%%%%%%%%%%%%%%%%%%%%%%% 

%\subsection{Perfect Mixability of Near-Final Configurations}
%\label{subsec: near-final configurations}
%
%%%%%%%%%%%%%%%%%%%%%%%%%%%%%%%%%%%%%%%%%%%%%%%%%%%%%%%%%%%%%%%%%%%%%%%%%%%%%%
%%%%%%%%%%%%%%%%%%%%%%%%%%%%%%%%%%%%%%%%%%%%%%%%%%%%%%%%%%%%%%%%%%%%%%%%%%%%%% 

%\subsection{Perfect Mixability of Near-Final Configurations}
%\label{subsec: near-final configurations}
%
%%%%%%%%%%%%%%%%%%%%%%%%%%%%%%%%%%%%%%%%%%%%%%%%%%%%%%%%%%%%%%%%%%%%%%%%%%%%%%
%%%%%%%%%%%%%%%%%%%%%%%%%%%%%%%%%%%%%%%%%%%%%%%%%%%%%%%%%%%%%%%%%%%%%%%%%%%%%% 

%\subsection{Perfect Mixability of Near-Final Configurations}
%\label{subsec: near-final configurations}
%\input{5-1_near-final_configurations.tex}

Let $C\ingroundset\integers$ be a configuration with
$\nofdroplets{C} = n = \sigma 2^\tau$, for some odd $\sigma\in\posintegers$ and $\tau\in\nonnegintegers$,
with $\average(C) = \mu\in\integers$.
We say that $C$ is \emph{near-final} if it can be partitioned
into disjoint multisets $C_1,C_2,...,C_k$, such that, for each $j$,
$\average(C_j) = \mu$ and $\nofdroplets{C_j}$ is a power of $2$.
In this sub-section we show (Lemma~\ref{lem: mixability for near-final} below)
that near-final configurations are perfectly mixable with precision $0$.
We also show that configurations with
only two different concentrations that satisfy Condition~{\MixCond} are near-final,
and thus perfectly mixable.

Define $\potential(C) = \sum_{c\in C}(c - \mu)^2$, which can be thought of as the
un-normalized variance of $C$. Obviously $\potential(C) \in \nonnegintegers$,
$\potential(C)=0$ if and only if $C$ is a perfect mixture,
and, by a straightforward calculation, mixing any two different same-parity concentrations in $C$ 
decreases the value of $\potential(C)$ by at least $1$.

%%%%%%%%%%%

\begin{lemma}\label{lem: mixability for near-final} 
	If $C$ is near-final
	then $C$ is perfectly mixable with precision $0$.
\end{lemma}

\begin{proof}
It is sufficient to prove the lemma for the case when $n$ is a power
of $2$. (Otherwise, we can apply it separately to each set $C_j$ in the 
partition of $C$ from the definition of near-final configurations.)	

So assume that $n$ is a power of $2$.
It is sufficient to show that if $\nofconcentrations{C} = m \neq 1$ (that is,
$C$ is not yet perfectly mixed) then $C$ contains
	two different concentrations with the same parity. 
	(Each such mixing strictly decreases $\potential(C)$, so a finite sequence
	of such mixing operations will perfectly mix $C$.)
	This is trivially true when $m\geq 3$,
	so it is sufficient to prove it for $m=2$,
	that is for $C = \braced{f_1:c_1, f_2:c_2}$. Without loss of generality, 
	by Observation~\ref{obs: offsetting does not affect reachability},
	we can assume that $c_2 = 0$,
	and then we claim that $c_1$ is even.
	We have $\average(C) = \mu = f_1 c_1/n$.
	As $\mu\in\integers$, $f_1 < n$ and $n$ is a power of $2$,
	we have that $c_1$ must be even, as claimed.
\end{proof}

%%%%%%%%%%

\begin{lemma}\label{lem: mixability m = 2}
	Assume that $\nofconcentrations{C}=2$, say $C=\braced{f_1:c_1,f_2:c_2}$,
	and that $C$ satisfies Condition~{\MixCond}.
	Then $\sigma$ divides $f_1$ and $f_2$.
	Consequently, we have that
	$n$ is not prime and $C$ is near-final.
\end{lemma}

\begin{proof}
Without loss of generality, assume that $c_2 = 0$. (Otherwise consider
$C' = C-c_2$ instead. This does not affect Condition~~{\MixCond} and the property of being near-final.)
Let $c_1 = \alpha 2^{\gamma}$, for some odd $\alpha\in\posintegers$. 
Then $\mu = f_1c_1/n = f_1\alpha 2^{\gamma}/(\sigma 2^\tau)$. 
Since $\alpha$ divides $c_1$ and $c_2$, Condition~{\MixCond} implies that 
$\alpha$ must also divide $\mu$. In other words, $\mu/\alpha = f_1 2^{\gamma}/(\sigma 2^\tau)$
is integer. This implies, in turn, that $f_1$ is a multiple of $\sigma$, as claimed.
Since $f_2 = n-f_1$, it also gives us that $f_2$ is a multiple of $\sigma$.    

This immediately implies that $n$ cannot be prime, for $n=f_1+f_2$ is a sum of
two non-zero multiples of $\sigma$.

To prove the last claim,
let $f_1 =\sigma f'_1$ and $f_2 = \sigma f'_2$,  for some $f'_1, f'_2 \in\posintegers$.
Partition $C$ into $\sigma$ sub-multisets of the form $C_j = \braced{f_1':c_1,f'_2:0}$, for $j=1,2,...,\sigma$.
The cardinality of each $C_j$ is $f'_1+f'_2 = n/\sigma = 2^\tau$ and
its average is $\average(C_j) = f'_1c_1/(f'_1+f'_2) = (f_1/\sigma)c_1/(n/\sigma) = f_1c_1/n = \mu$.
Therefore $C$ is near-final, as claimed.
\end{proof}

We remark that when $n$ is a power of $2$ there is an alternative way to
perfectly mix $C$, using a divide-and-conquer approach: partition $C$ into
two equal-size multisets $C', C''$, mix each of these recursively, obtaining
$n/2$ droplets with concentration $\mu' = \average(C')$ and $n/2$ droplets
with concentration $\mu'' = \average(C'')$, 
and then mix $n/2$ disjoint pairs of droplets, one $\mu'$ and the other $\mu''$,
producing $n$ droplets with concentration $\mu = \average(C)$.
This approach, however, produces mixing graphs where the intermediate precision 
could be quite large, so it is not sufficient for our purpose.






Let $C\ingroundset\integers$ be a configuration with
$\nofdroplets{C} = n = \sigma 2^\tau$, for some odd $\sigma\in\posintegers$ and $\tau\in\nonnegintegers$,
with $\average(C) = \mu\in\integers$.
We say that $C$ is \emph{near-final} if it can be partitioned
into disjoint multisets $C_1,C_2,...,C_k$, such that, for each $j$,
$\average(C_j) = \mu$ and $\nofdroplets{C_j}$ is a power of $2$.
In this sub-section we show (Lemma~\ref{lem: mixability for near-final} below)
that near-final configurations are perfectly mixable with precision $0$.
We also show that configurations with
only two different concentrations that satisfy Condition~{\MixCond} are near-final,
and thus perfectly mixable.

Define $\potential(C) = \sum_{c\in C}(c - \mu)^2$, which can be thought of as the
un-normalized variance of $C$. Obviously $\potential(C) \in \nonnegintegers$,
$\potential(C)=0$ if and only if $C$ is a perfect mixture,
and, by a straightforward calculation, mixing any two different same-parity concentrations in $C$ 
decreases the value of $\potential(C)$ by at least $1$.

%%%%%%%%%%%

\begin{lemma}\label{lem: mixability for near-final} 
	If $C$ is near-final
	then $C$ is perfectly mixable with precision $0$.
\end{lemma}

\begin{proof}
It is sufficient to prove the lemma for the case when $n$ is a power
of $2$. (Otherwise, we can apply it separately to each set $C_j$ in the 
partition of $C$ from the definition of near-final configurations.)	

So assume that $n$ is a power of $2$.
It is sufficient to show that if $\nofconcentrations{C} = m \neq 1$ (that is,
$C$ is not yet perfectly mixed) then $C$ contains
	two different concentrations with the same parity. 
	(Each such mixing strictly decreases $\potential(C)$, so a finite sequence
	of such mixing operations will perfectly mix $C$.)
	This is trivially true when $m\geq 3$,
	so it is sufficient to prove it for $m=2$,
	that is for $C = \braced{f_1:c_1, f_2:c_2}$. Without loss of generality, 
	by Observation~\ref{obs: offsetting does not affect reachability},
	we can assume that $c_2 = 0$,
	and then we claim that $c_1$ is even.
	We have $\average(C) = \mu = f_1 c_1/n$.
	As $\mu\in\integers$, $f_1 < n$ and $n$ is a power of $2$,
	we have that $c_1$ must be even, as claimed.
\end{proof}

%%%%%%%%%%

\begin{lemma}\label{lem: mixability m = 2}
	Assume that $\nofconcentrations{C}=2$, say $C=\braced{f_1:c_1,f_2:c_2}$,
	and that $C$ satisfies Condition~{\MixCond}.
	Then $\sigma$ divides $f_1$ and $f_2$.
	Consequently, we have that
	$n$ is not prime and $C$ is near-final.
\end{lemma}

\begin{proof}
Without loss of generality, assume that $c_2 = 0$. (Otherwise consider
$C' = C-c_2$ instead. This does not affect Condition~~{\MixCond} and the property of being near-final.)
Let $c_1 = \alpha 2^{\gamma}$, for some odd $\alpha\in\posintegers$. 
Then $\mu = f_1c_1/n = f_1\alpha 2^{\gamma}/(\sigma 2^\tau)$. 
Since $\alpha$ divides $c_1$ and $c_2$, Condition~{\MixCond} implies that 
$\alpha$ must also divide $\mu$. In other words, $\mu/\alpha = f_1 2^{\gamma}/(\sigma 2^\tau)$
is integer. This implies, in turn, that $f_1$ is a multiple of $\sigma$, as claimed.
Since $f_2 = n-f_1$, it also gives us that $f_2$ is a multiple of $\sigma$.    

This immediately implies that $n$ cannot be prime, for $n=f_1+f_2$ is a sum of
two non-zero multiples of $\sigma$.

To prove the last claim,
let $f_1 =\sigma f'_1$ and $f_2 = \sigma f'_2$,  for some $f'_1, f'_2 \in\posintegers$.
Partition $C$ into $\sigma$ sub-multisets of the form $C_j = \braced{f_1':c_1,f'_2:0}$, for $j=1,2,...,\sigma$.
The cardinality of each $C_j$ is $f'_1+f'_2 = n/\sigma = 2^\tau$ and
its average is $\average(C_j) = f'_1c_1/(f'_1+f'_2) = (f_1/\sigma)c_1/(n/\sigma) = f_1c_1/n = \mu$.
Therefore $C$ is near-final, as claimed.
\end{proof}

We remark that when $n$ is a power of $2$ there is an alternative way to
perfectly mix $C$, using a divide-and-conquer approach: partition $C$ into
two equal-size multisets $C', C''$, mix each of these recursively, obtaining
$n/2$ droplets with concentration $\mu' = \average(C')$ and $n/2$ droplets
with concentration $\mu'' = \average(C'')$, 
and then mix $n/2$ disjoint pairs of droplets, one $\mu'$ and the other $\mu''$,
producing $n$ droplets with concentration $\mu = \average(C)$.
This approach, however, produces mixing graphs where the intermediate precision 
could be quite large, so it is not sufficient for our purpose.






Let $C\ingroundset\integers$ be a configuration with
$\nofdroplets{C} = n = \sigma 2^\tau$, for some odd $\sigma\in\posintegers$ and $\tau\in\nonnegintegers$,
with $\average(C) = \mu\in\integers$.
We say that $C$ is \emph{near-final} if it can be partitioned
into disjoint multisets $C_1,C_2,...,C_k$, such that, for each $j$,
$\average(C_j) = \mu$ and $\nofdroplets{C_j}$ is a power of $2$.
In this sub-section we show (Lemma~\ref{lem: mixability for near-final} below)
that near-final configurations are perfectly mixable with precision $0$.
We also show that configurations with
only two different concentrations that satisfy Condition~{\MixCond} are near-final,
and thus perfectly mixable.

Define $\potential(C) = \sum_{c\in C}(c - \mu)^2$, which can be thought of as the
un-normalized variance of $C$. Obviously $\potential(C) \in \nonnegintegers$,
$\potential(C)=0$ if and only if $C$ is a perfect mixture,
and, by a straightforward calculation, mixing any two different same-parity concentrations in $C$ 
decreases the value of $\potential(C)$ by at least $1$.

%%%%%%%%%%%

\begin{lemma}\label{lem: mixability for near-final} 
	If $C$ is near-final
	then $C$ is perfectly mixable with precision $0$.
\end{lemma}

\begin{proof}
It is sufficient to prove the lemma for the case when $n$ is a power
of $2$. (Otherwise, we can apply it separately to each set $C_j$ in the 
partition of $C$ from the definition of near-final configurations.)	

So assume that $n$ is a power of $2$.
It is sufficient to show that if $\nofconcentrations{C} = m \neq 1$ (that is,
$C$ is not yet perfectly mixed) then $C$ contains
	two different concentrations with the same parity. 
	(Each such mixing strictly decreases $\potential(C)$, so a finite sequence
	of such mixing operations will perfectly mix $C$.)
	This is trivially true when $m\geq 3$,
	so it is sufficient to prove it for $m=2$,
	that is for $C = \braced{f_1:c_1, f_2:c_2}$. Without loss of generality, 
	by Observation~\ref{obs: offsetting does not affect reachability},
	we can assume that $c_2 = 0$,
	and then we claim that $c_1$ is even.
	We have $\average(C) = \mu = f_1 c_1/n$.
	As $\mu\in\integers$, $f_1 < n$ and $n$ is a power of $2$,
	we have that $c_1$ must be even, as claimed.
\end{proof}

%%%%%%%%%%

\begin{lemma}\label{lem: mixability m = 2}
	Assume that $\nofconcentrations{C}=2$, say $C=\braced{f_1:c_1,f_2:c_2}$,
	and that $C$ satisfies Condition~{\MixCond}.
	Then $\sigma$ divides $f_1$ and $f_2$.
	Consequently, we have that
	$n$ is not prime and $C$ is near-final.
\end{lemma}

\begin{proof}
Without loss of generality, assume that $c_2 = 0$. (Otherwise consider
$C' = C-c_2$ instead. This does not affect Condition~~{\MixCond} and the property of being near-final.)
Let $c_1 = \alpha 2^{\gamma}$, for some odd $\alpha\in\posintegers$. 
Then $\mu = f_1c_1/n = f_1\alpha 2^{\gamma}/(\sigma 2^\tau)$. 
Since $\alpha$ divides $c_1$ and $c_2$, Condition~{\MixCond} implies that 
$\alpha$ must also divide $\mu$. In other words, $\mu/\alpha = f_1 2^{\gamma}/(\sigma 2^\tau)$
is integer. This implies, in turn, that $f_1$ is a multiple of $\sigma$, as claimed.
Since $f_2 = n-f_1$, it also gives us that $f_2$ is a multiple of $\sigma$.    

This immediately implies that $n$ cannot be prime, for $n=f_1+f_2$ is a sum of
two non-zero multiples of $\sigma$.

To prove the last claim,
let $f_1 =\sigma f'_1$ and $f_2 = \sigma f'_2$,  for some $f'_1, f'_2 \in\posintegers$.
Partition $C$ into $\sigma$ sub-multisets of the form $C_j = \braced{f_1':c_1,f'_2:0}$, for $j=1,2,...,\sigma$.
The cardinality of each $C_j$ is $f'_1+f'_2 = n/\sigma = 2^\tau$ and
its average is $\average(C_j) = f'_1c_1/(f'_1+f'_2) = (f_1/\sigma)c_1/(n/\sigma) = f_1c_1/n = \mu$.
Therefore $C$ is near-final, as claimed.
\end{proof}

We remark that when $n$ is a power of $2$ there is an alternative way to
perfectly mix $C$, using a divide-and-conquer approach: partition $C$ into
two equal-size multisets $C', C''$, mix each of these recursively, obtaining
$n/2$ droplets with concentration $\mu' = \average(C')$ and $n/2$ droplets
with concentration $\mu'' = \average(C'')$, 
and then mix $n/2$ disjoint pairs of droplets, one $\mu'$ and the other $\mu''$,
producing $n$ droplets with concentration $\mu = \average(C)$.
This approach, however, produces mixing graphs where the intermediate precision 
could be quite large, so it is not sufficient for our purpose.






Let $C\ingroundset\integers$ be a configuration with
$\nofdroplets{C} = n = \sigma 2^\tau$, for some odd $\sigma\in\posintegers$ and $\tau\in\nonnegintegers$,
with $\average(C) = \mu\in\integers$.
We say that $C$ is \emph{near-final} if it can be partitioned
into disjoint multisets $C_1,C_2,...,C_k$, such that, for each $j$,
$\average(C_j) = \mu$ and $\nofdroplets{C_j}$ is a power of $2$.
In this sub-section we show (Lemma~\ref{lem: mixability for near-final} below)
that near-final configurations are perfectly mixable with precision $0$.
We also show that configurations with
only two different concentrations that satisfy Condition~{\MixCond} are near-final,
and thus perfectly mixable.

Define $\potential(C) = \sum_{c\in C}(c - \mu)^2$, which can be thought of as the
un-normalized variance of $C$. Obviously $\potential(C) \in \nonnegintegers$,
$\potential(C)=0$ if and only if $C$ is a perfect mixture,
and, by a straightforward calculation, mixing any two different same-parity concentrations in $C$ 
decreases the value of $\potential(C)$ by at least $1$.

%%%%%%%%%%%

\begin{lemma}\label{lem: mixability for near-final} 
	If $C$ is near-final
	then $C$ is perfectly mixable with precision $0$.
\end{lemma}

\begin{proof}
It is sufficient to prove the lemma for the case when $n$ is a power
of $2$. (Otherwise, we can apply it separately to each set $C_j$ in the 
partition of $C$ from the definition of near-final configurations.)	

So assume that $n$ is a power of $2$.
It is sufficient to show that if $\nofconcentrations{C} = m \neq 1$ (that is,
$C$ is not yet perfectly mixed) then $C$ contains
	two different concentrations with the same parity. 
	(Each such mixing strictly decreases $\potential(C)$, so a finite sequence
	of such mixing operations will perfectly mix $C$.)
	This is trivially true when $m\geq 3$,
	so it is sufficient to prove it for $m=2$,
	that is for $C = \braced{f_1:c_1, f_2:c_2}$. Without loss of generality, 
	by Observation~\ref{obs: offsetting does not affect reachability},
	we can assume that $c_2 = 0$,
	and then we claim that $c_1$ is even.
	We have $\average(C) = \mu = f_1 c_1/n$.
	As $\mu\in\integers$, $f_1 < n$ and $n$ is a power of $2$,
	we have that $c_1$ must be even, as claimed.
\end{proof}

%%%%%%%%%%

\begin{lemma}\label{lem: mixability m = 2}
	Assume that $\nofconcentrations{C}=2$, say $C=\braced{f_1:c_1,f_2:c_2}$,
	and that $C$ satisfies Condition~{\MixCond}.
	Then $\sigma$ divides $f_1$ and $f_2$.
	Consequently, we have that
	$n$ is not prime and $C$ is near-final.
\end{lemma}

\begin{proof}
Without loss of generality, assume that $c_2 = 0$. (Otherwise consider
$C' = C-c_2$ instead. This does not affect Condition~~{\MixCond} and the property of being near-final.)
Let $c_1 = \alpha 2^{\gamma}$, for some odd $\alpha\in\posintegers$. 
Then $\mu = f_1c_1/n = f_1\alpha 2^{\gamma}/(\sigma 2^\tau)$. 
Since $\alpha$ divides $c_1$ and $c_2$, Condition~{\MixCond} implies that 
$\alpha$ must also divide $\mu$. In other words, $\mu/\alpha = f_1 2^{\gamma}/(\sigma 2^\tau)$
is integer. This implies, in turn, that $f_1$ is a multiple of $\sigma$, as claimed.
Since $f_2 = n-f_1$, it also gives us that $f_2$ is a multiple of $\sigma$.    

This immediately implies that $n$ cannot be prime, for $n=f_1+f_2$ is a sum of
two non-zero multiples of $\sigma$.

To prove the last claim,
let $f_1 =\sigma f'_1$ and $f_2 = \sigma f'_2$,  for some $f'_1, f'_2 \in\posintegers$.
Partition $C$ into $\sigma$ sub-multisets of the form $C_j = \braced{f_1':c_1,f'_2:0}$, for $j=1,2,...,\sigma$.
The cardinality of each $C_j$ is $f'_1+f'_2 = n/\sigma = 2^\tau$ and
its average is $\average(C_j) = f'_1c_1/(f'_1+f'_2) = (f_1/\sigma)c_1/(n/\sigma) = f_1c_1/n = \mu$.
Therefore $C$ is near-final, as claimed.
\end{proof}

We remark that when $n$ is a power of $2$ there is an alternative way to
perfectly mix $C$, using a divide-and-conquer approach: partition $C$ into
two equal-size multisets $C', C''$, mix each of these recursively, obtaining
$n/2$ droplets with concentration $\mu' = \average(C')$ and $n/2$ droplets
with concentration $\mu'' = \average(C'')$, 
and then mix $n/2$ disjoint pairs of droplets, one $\mu'$ and the other $\mu''$,
producing $n$ droplets with concentration $\mu = \average(C)$.
This approach, however, produces mixing graphs where the intermediate precision 
could be quite large, so it is not sufficient for our purpose.




